\chapter{Introducción}
\label{cap:introduccion}

\section{Contexto}
Los sistemas de composición tipográfica basados en \TeX{} siguen siendo la
herramienta de referencia para publicar artículos científicos, tesis y
documentación técnica \citep{lamport1994}. Sin embargo, instalar y mantener
una distribución completa de \TeX{} Live exige varios gigabytes de disco y un
conocimiento considerable de su ecosistema de paquetes. Para una organización
que solo necesita transformar una plantilla y unos datos en un archivo PDF
(por ejemplo, un sistema de generación de certificados o de informes
periódicos), esa carga resulta desproporcionada.

Durante los últimos años han aparecido editores en línea que centralizan la
compilación \citep{overleaf2023}, pero sus componentes internos no se exponen
como piezas reutilizables. Este informe estudia la compilación como un
servicio independiente, sin estado, que recibe un conjunto de archivos y
devuelve el PDF resultante junto con un diagnóstico estructurado de errores y
advertencias.

\section{Objetivos}
El trabajo persigue tres objetivos concretos:

\begin{enumerate}[label=\textbf{O\arabic*.}]
  \item Diseñar un modelo de ejecución que permita compilar documentos de
    origen no confiable en una infraestructura compartida sin comprometer el
    servidor ni los datos de otros usuarios.
  \item Garantizar que una misma entrada produzca siempre la misma salida
    mientras no cambie la versión de la imagen de compilación.
  \item Reducir la latencia percibida mediante una caché de artefactos que
    evite recompilar documentos que no han cambiado.
\end{enumerate}

\section{Alcance y estructura del informe}
El alcance se limita a los tres motores más utilizados (pdf\TeX{}, Xe\TeX{} y
Lua\TeX{}) y al flujo de trabajo gestionado por \texttt{latexmk}. Quedan fuera
del estudio las funcionalidades que requieren ejecutar programas externos
durante la compilación, como \texttt{minted} o la externalización de figuras
de TikZ, porque el modelo de aislamiento las prohíbe de forma deliberada.

El resto del documento se organiza de la siguiente manera. El
capítulo~\ref{cap:diseno} describe la arquitectura del servicio, el modelo de
aislamiento y la clave de caché. El capítulo~\ref{cap:resultados} presenta la
metodología de evaluación, los resultados obtenidos y las lecciones
aprendidas durante seis meses de operación.
